%% 16장 자료구조
\chapter{자료구조}
\label{ch:data-structures}

\abstract{알고리즘의 속도는 데이터를 어떤 구조에 담느냐에 크게 좌우된다. 이 장에서는 수치 라이브러리와 데이터프레임 라이브러리가 가장 많이 쓰는 세 자료구조를 다룬다. 첫째는 가장 작은 원소를 빠르게 꺼내는 힙이다. 배열 하나로 완전 이진 나무를 나타내는 방법, 힙을 선형 시간에 만드는 이유, 원소의 키를 줄이는 연산을 지원하는 색인 힙을 본다. 둘째는 키로 값을 찾는 해시 테이블이다. 해시 함수가 갖추어야 할 성질과 열린 주소법의 분석을 다룬 뒤, 분기 예측 실패를 피하려고 SIMD로 여러 칸을 한꺼번에 비교하는 Swiss table을 살핀다. 셋째는 원소들을 서로소인 집합으로 나누고 합치는 합집합--찾기이다. 각 구조를 stardust의 데이터프레임 그룹화와 그래프 알고리즘에서 실제로 어떻게 쓰는지와 함께 본다.}

\section{힙}
\label{sec:heap}

\subsection{배열 속의 나무}

원소들을 넣고 그 가운데 가장 작은 것을 꺼내는 일을 되풀이하는 자료구조를 \term{우선순위 큐}{priority queue}라 한다. 정렬된 배열은 꺼내기가 빠르지만 넣기가 $O(n)$이고, 정렬하지 않은 배열은 그 반대이다. \term{이진 힙}{binary heap}은 두 연산을 모두 $O(\log n)$에 한다.

이진 힙은 \term{힙 성질}{heap property}, 곧 모든 마디의 키가 자식의 키보다 크지 않다는 성질을 만족하는 완전 이진 나무이다. 완전 이진 나무는 마지막 층을 빼면 꽉 차 있고 마지막 층은 왼쪽부터 채워져 있으므로, 층 순서로 배열에 담으면 포인터가 필요 없다. 첨자를 $0$부터 매기면 마디 $i$의 자식은 $2i+1$과 $2i+2$이고 부모는 $\lfloor(i-1)/2\rfloor$이다. 힙 성질에서 가장 작은 원소는 배열의 첫 칸인 뿌리에 있다.

새 원소는 배열의 끝에 넣은 뒤 부모보다 작은 동안 부모와 자리를 바꾸며 위로 올린다. 뿌리를 꺼낼 때는 마지막 원소를 뿌리 자리로 옮긴 뒤 더 작은 자식보다 큰 동안 그 자식과 자리를 바꾸며 아래로 내린다. 나무의 높이가 $\lfloor\log_2n\rfloor$이므로 두 연산 모두 $O(\log n)$이다.

\subsection{선형 시간에 힙 만들기}

$n$개의 원소로 힙을 만들 때 하나씩 넣으면 $O(n\log n)$이지만, 배열을 그대로 두고 아래쪽 마디부터 차례로 ``아래로 내리기''를 하면 선형 시간이면 된다\cite{floyd1964treesort}.

\begin{proposition}
\label{prop:heapify}
잎이 아닌 마디를 마지막부터 뿌리까지 차례로 아래로 내리면 힙이 만들어지고, 비교 횟수는 $O(n)$이다.
\end{proposition}

\begin{proof}
마디 $i$를 내릴 때 그 두 자식을 뿌리로 하는 부분 나무는 이미 힙이므로, 내리기가 끝나면 $i$를 뿌리로 하는 부분 나무도 힙이 된다. 따라서 뿌리까지 마치면 전체가 힙이다. 높이가 $h$인 마디를 내리는 데는 많아야 $2h$번 비교하고, 높이가 $h$인 마디는 많아야 $\lceil n/2^{h+1}\rceil$개이다. 그러므로 비교 횟수는
\[
\sum_{h\ge1}\Bigl\lceil\frac n{2^{h+1}}\Bigr\rceil\cdot2h\le\sum_{h\ge1}\Bigl(\frac n{2^{h+1}}+1\Bigr)2h=n\sum_{h\ge1}\frac h{2^h}+O(\log^2n)=2n+O(\log^2n)
\]
이다. 마지막 등호는 $\sum_{h\ge1}h x^h=x/(1-x)^2$에 $x=1/2$을 넣은 것이다.
\end{proof}

힙을 만든 뒤 뿌리를 꺼내 배열의 끝으로 보내는 일을 $n$번 되풀이하면 \term{힙 정렬}{heapsort}이 되고\cite{williams1964heapsort}, 추가 메모리 없이 최악의 경우에도 $O(n\log n)$이다. 15장의 인트로 정렬이 최악의 경우를 막는 데 쓰는 것이 이것이다. 다만 힙 정렬은 배열의 먼 곳을 오가므로 캐시에 불리해서 평균적으로는 퀵정렬보다 느리다.

\subsection{$d$진 힙과 색인 힙}

자식이 $d$개인 \term{$d$진 힙}{$d$-ary heap}에서는 나무의 높이가 $\log_dn$으로 낮아져 위로 올리기가 빨라지는 대신, 아래로 내릴 때 자식 $d$개 가운데 가장 작은 것을 찾아야 한다\cite{johnson1975dheap}. $d=4$이면 높이가 절반이 되고, 자식 네 개의 키가 배열에서 이웃해 있으므로 binary64 키라면 $32$바이트, 곧 캐시 줄 하나 안에 든다. 13장에서 보았듯이 캐시를 놓치는 비용이 비교의 비용보다 훨씬 크므로, 비교가 늘어도 메모리 접근이 줄어 실제로는 더 빠른 경우가 많다.

최단 경로를 찾는 Dijkstra의 알고리즘은 우선순위 큐에 마디를 거리 순으로 넣고 꺼내는데\cite{dijkstra1959}, 어떤 마디의 거리가 더 짧게 갱신되면 큐에 이미 있는 그 마디의 키를 줄여야 한다. 키를 줄이려면 그 원소가 배열의 어디에 있는지 알아야 하므로, 마디 번호에서 힙 안의 자리로 가는 배열 \code{pos}를 함께 둔다. 자리를 바꿀 때마다 \code{pos}도 함께 고치면, 키 줄이기는 \code{pos}로 자리를 찾은 뒤 위로 올리기 한 번이다. 이것을 \term{색인 힙}{indexed heap}이라 한다. 키 줄이기를 상수 분할상환 시간에 하는 Fibonacci 힙\cite{fredman1987fibonacci}이 이론적으로는 더 좋지만, 상수 인자와 메모리 접근 때문에 실제로는 배열 기반 힙이 거의 언제나 빠르다.

\begin{instardust}[\code{csrc/graph/paths.hpp}]
그래프 영역의 Dijkstra는 4진 색인 힙을 쓴다. 키 배열 \code{hk}와 마디 배열 \code{hv}를 따로 두고, \code{pos[v]}에 마디 $v$의 자리를, 큐에 없으면 $-1$을 적는다. 마디 하나는 큐에 많아야 한 번 들어가고, 더 작은 키가 오면 그 자리에서 위로 올린다. 파일의 주석에는 이 장의 주제를 넘어서는 흥미로운 관찰이 있다. 거리는 반올림한 합 $d(v)=\min_u\fl(d(u)+w(u,v))$의 가장 작은 고정점으로 정의되는데, 이완할 간선이 하나도 남지 않을 때 멈추는 이완 알고리즘은 이완의 순서와 상관없이 모두 이 고정점에서 멈춘다. 그래서 색인 힙 Dijkstra, 대기열 Bellman--Ford, 힙 없는 조밀 Dijkstra가 모두 같은 비트의 거리를 낸다. 14장의 결정성은 알고리즘의 선택과도 무관하게 성립한다.
\end{instardust}

\section{해시 테이블}
\label{sec:hash-table}

\subsection{해시 함수}

\term{해시 테이블}{hash table}은 키를 \term{해시 함수}{hash function} $h$로 정수로 바꾸어 그 값으로 배열의 칸을 골라 키와 값을 담는다. 좋은 해시 함수라면 서로 다른 키들이 칸들에 고르게 퍼지므로 넣기와 찾기가 평균 $O(1)$이다. 두 키가 같은 칸으로 가는 \term{충돌}{collision}은 피할 수 없으므로 그것을 처리하는 방법이 필요하다.

해시 함수에 요구하는 성질은 쓰임에 따라 다르다. 공격자가 충돌을 일부러 만들 수 있는 환경에서는 무작위로 고른 함수족을 쓰는 \term{유니버설 해싱}{universal hashing}이 이론적 보장을 준다\cite{carter1979universal}. 데이터 분석처럼 그런 걱정이 덜한 곳에서는 빠르면서도 키의 패턴이 출력에 드러나지 않는 함수가 필요하다. 여기서 중요한 성질이 \term{쇄도 효과}{avalanche effect}로, 입력의 비트 하나를 바꾸면 출력의 각 비트가 각각 반의 확률로 바뀌는 것이다. 이 성질이 있으면 $0,1,2,\dots$처럼 연속된 정수 키도 출력의 아래쪽 비트와 위쪽 비트에 고르게 퍼진다.

2장에서 본 MurmurHash3의 마무리 함수 fmix64는 곱셈 두 번과 이동--XOR 세 번으로 된 전단사이다\cite{appleby-smhasher}. 무작위 $64$비트 입력 $2$만 개에서 비트 하나씩을 바꾸어 보면 출력 비트가 평균 $32.0$개 바뀌어 쇄도 효과가 거의 완벽하고, 입력 $0,1,2,\dots,7$의 출력의 아래 $7$비트는 $0,44,103,78,117,117,115,93$으로 흩어진다. 문자열처럼 길이가 정해지지 않은 키는 $8$바이트씩 읽어 섞는데, 두 $64$비트 값을 곱한 $128$비트 결과의 위쪽과 아래쪽을 XOR하는 것이 빠르고 섞임이 좋은 방법으로 알려져 있다.

\subsection{열린 주소법}

충돌을 처리하는 방법은 크게 둘이다. 칸마다 같은 칸으로 온 키들의 목록을 두는 \term{체인법}{chaining}과, 모든 키를 배열 자체에 넣되 칸이 차 있으면 정해진 규칙으로 다른 칸을 찾아가는 \term{열린 주소법}{open addressing}이다\cite{peterson1957}. 열린 주소법은 포인터가 없어 메모리를 덜 쓰고, 이웃한 칸을 차례로 보는 \term{선형 탐사}{linear probing}를 쓰면 탐색이 캐시 줄 안에서 이루어져 빠르다.

선형 탐사의 비용은 테이블이 찬 비율, 곧 \term{적재율}{load factor} $\alpha$에 달려 있다. 해시 값이 균등하고 독립이라는 가정에서 Knuth는 찾는 키가 있을 때와 없을 때의 평균 탐사 횟수가 각각 대략
\[
\frac12\Bigl(1+\frac1{1-\alpha}\Bigr),\qquad\frac12\Bigl(1+\frac1{(1-\alpha)^2}\Bigr)
\]
임을 보였다\cite{knuth1998taocp3}. $\alpha=0.5$이면 $1.5$번과 $2.5$번이지만 $\alpha=0.9$이면 $5.5$번과 $50.5$번으로, 테이블이 차면 탐사가 급격히 늘어난다. 차 있는 칸들이 덩어리를 이루어 그 덩어리가 더 빨리 커지는 \term{군집}{clustering} 때문이다. 그래서 선형 탐사 테이블은 대개 적재율이 $1/2$ 정도를 넘으면 크기를 두 배로 늘린다.

\subsection{Swiss table}

탐사 횟수만 보면 선형 탐사는 적재율이 낮을 때 거의 최적이다. 그런데 실제 성능은 탐사 횟수가 아니라 13장에서 본 분기 예측에 좌우될 수 있다. 선형 탐사는 칸마다 ``이 칸이 비었나'', ``키가 같나''를 분기로 묻는데, 첫 칸에서 찾을지 다음 칸으로 갈지는 예측하기 어렵다.

Google의 C++ 라이브러리 Abseil에서 처음 쓰인 \term{Swiss table}{Swiss table}은 이 분기를 SIMD로 없앤다. 칸마다 $1$바이트의 \term{제어 바이트}{control byte}를 따로 두는데, 빈 칸이면 $0$이고 찬 칸이면 그 키의 해시 값의 아래 $7$비트에 맨 위 비트를 세운 값, 곧 \term{태그}{tag}이다. 칸 $16$개를 한 \term{그룹}{group}으로 묶고 해시 값의 나머지 비트로 그룹을 고른다. 키를 찾을 때는 그룹의 제어 바이트 $16$개를 레지스터 하나에 읽어 찾는 키의 태그와 한 번에 비교하고, 일치한 칸에서만 실제 키를 비교한다. 태그 $7$비트가 우연히 같을 확률은 $1/128$이므로 실제 키 비교는 거의 일치하는 칸 하나에서만 일어난다. 그룹에 빈 칸이 하나라도 있으면 키가 테이블에 없다는 뜻이므로 거기서 멈추고, 없으면 다음 그룹으로 간다. 남는 분기는 ``후보가 있나''와 ``빈 칸이 있나''뿐이고, 둘 다 거의 언제나 같은 쪽으로 가므로 잘 예측된다.

\begin{instardust}[\code{docs/design/frame.md} §8.3]
데이터프레임의 그룹화, 중복 찾기, 병합은 모두 하나의 \code{factorize}, 곧 키를 처음 나타난 순서대로 $0,1,2,\dots$의 정수 코드로 바꾸는 연산 위에 서 있고, 그 테이블이 Swiss table이다. 설계 문서에는 Swiss table을 고른 이유가 측정으로 남아 있다. 선형 탐사 테이블은 행 백만 개, 키 천 종에서 테이블이 L1 캐시에 다 들어가는데도 행마다 약 $19$사이클이 걸렸는데, 평균 탐사 횟수는 $1.48$로 적재율 $0.49$에서 위의 공식이 주는 값 그대로였다. 프로파일을 보니 시간의 $2/3$가 ``첫 칸의 키가 다르다'' 분기 바로 뒤에 몰려 있었고, Swiss table로 바꾸자 약 $2$배 빨라졌다. NEON에는 x86의 \code{movemask}처럼 비교 결과를 비트열로 모으는 명령이 없어서, 바이트 $16$개의 비교 결과를 \code{vshrn\_n\_u16}으로 칸마다 $4$비트씩 $64$비트 정수 하나로 줄인 뒤 \code{ctz}로 다음 후보를 찾는다. 빈 칸의 제어 바이트가 $0$이므로 운영체제가 $0$으로 채워 주는 새 페이지가 곧 빈 테이블이어서 크게 잡아 두어도 실제로 쓴 페이지만 메모리를 차지한다. 해시 값의 위쪽 $8$비트는 병렬 처리를 위해 키를 $256$개 구역으로 나누는 데, 아래 $7$비트는 태그에, 그 사이의 비트는 그룹을 고르는 데 쓰여 서로 겹치지 않는다.
\end{instardust}

\subsection{키의 정규화와 작은 범위의 정수}

해시 테이블은 ``비트가 같으면 같은 키''일 때 가장 단순하다. 그런데 부동소수점 키에서 \code{==}로 같은 두 값 $+0$과 $-0$은 비트가 다르고, NaN은 $2^{53}-2$가지의 비트 패턴이 있는데 데이터 분석에서는 보통 하나의 그룹으로 묶는다(8장). 그래서 해시하기 전에 키를 정규화한다. $-0$은 $+0$의 비트로, 모든 NaN은 정해진 NaN 하나의 비트로 바꾸면, 이후로는 비트 비교만으로 같은 키를 판정할 수 있다. 결과에는 원래 값이 필요하므로 각 코드가 처음 나온 행 번호를 함께 기억해 두었다가 그 행의 원래 값을 돌려준다.

키가 작은 범위의 정수이면 해시조차 필요 없다. 값이 $[lo,lo+s)$에 있으면 키 $lo+k$를 칸 $k$에 두는 \term{직접 주소 표}{direct-address table}가 해시 없이, 탐사 없이 행마다 메모리 접근 하나로 끝난다. 범위 $s$가 행의 수에 비해 너무 크지 않을 때만 쓸 수 있지만, 범주형 데이터의 코드나 작은 정수 키처럼 그런 경우가 실제로 흔하다.

\begin{instardust}[\code{docs/design/frame.md} §8.2, §8.9]
\code{float\_key}가 위의 정규화이다. NaN이면 \code{0x7ff8000000000000}을, $0.0$과 같으면 $0$을, 그 밖에는 비트 그대로를 돌려준다. int64와 bool 열의 값이 범위 $s\le\min(\max(2n,2^{16}),2^{28})$ 안에 있으면 직접 주소 표를 쓰는데, 칸에는 ``코드 더하기 $1$''을 $4$바이트로 담아 $0$을 ``아직 없음''으로 쓴다. 키가 많아 새 키가 나타나는지가 거의 무작위일 때는 ``새 키인가''의 분기마저 없애서, 모든 행이 칸과 다음 자리를 쓰되 새 키일 때만 끝 위치가 앞으로 가게 했더니 약 $2.6$배 빨라졌다고 설계 문서는 적는다. 어느 표를 쓰든 코드는 같으며, 처음 나타난 순서이므로 스레드 수와도 무관하다.
\end{instardust}

\section{합집합--찾기}
\label{sec:union-find}

원소 $0,\dots,n-1$을 서로소인 집합들로 나누어 두고, 두 원소가 같은 집합에 있는지 묻거나(찾기) 두 집합을 합치는(합집합) 연산을 되풀이하는 자료구조를 \term{합집합--찾기}{union--find} 또는 서로소 집합 자료구조라 한다\cite{galler1964}. 그래프의 연결 요소, 단일 연결 군집, 등가 관계의 계산에 쓰인다. 각 집합을 하나의 나무로 나타내고 원소마다 부모를 가리키는 배열 \code{parent}를 두며, 뿌리가 그 집합의 대표이다. 찾기는 부모를 따라 뿌리까지 올라가고, 합집합은 한 뿌리를 다른 뿌리의 자식으로 만든다.

아무렇게나 합치면 나무가 한 줄로 길어져 찾기가 $O(n)$이 된다. 두 가지 개선이 있다. 하나는 작은 나무를 큰 나무 밑에 붙이는 \term{크기에 따른 합치기}{union by size}로, 그러면 나무의 높이가 $O(\log n)$을 넘지 않는다. 원소가 더 깊은 쪽으로 옮겨질 때마다 그 원소가 속한 나무의 크기가 적어도 두 배가 되기 때문이다. 다른 하나는 찾기를 하면서 지나간 마디들을 뿌리에 더 가깝게 다시 붙이는 \term{경로 압축}{path compression}이다. 지나간 마디마다 부모를 할아버지로 바꾸는 \term{경로 반감}{path halving}이 가장 간단한 변형이다. Tarjan은 두 개선을 함께 쓰면 연산 $m$번의 총비용이 $O(m\,\alpha(n))$임을 보였는데\cite{tarjan1975,tarjan1984worst}, 여기서 $\alpha$는 Ackermann 함수의 역함수로서 우주의 원자 수만 한 $n$에서도 $4$를 넘지 않는다. 실용적으로는 상수 시간이다.

합집합--찾기의 결과, 곧 집합들의 분할 자체는 합치는 순서와 무관하다. 어떤 순서로 합치든 마지막 분할은 ``합친 쌍들이 만드는 그래프의 연결 요소''로 정해지기 때문이다. 그러나 각 집합의 대표인 뿌리가 어느 원소인지는 순서에 따라 달라질 수 있다. 결정성이 필요하면 대표를 분할만으로 정하면 되는데, 예컨대 언제나 집합의 가장 작은 원소를 뿌리로 삼으면 된다. 두 뿌리를 합칠 때 큰 뿌리를 작은 뿌리 밑에 붙이면 이것이 저절로 지켜진다. 이 규칙은 크기에 따른 합치기의 높이 보장을 포기하는 것이지만, 경로 압축이 있으면 실제 비용은 거의 차이가 없다.

\begin{instardust}[\code{csrc/graph/traversal.cpp}]
그래프 영역의 약한 연결 요소는 뿌리가 그 집합의 가장 작은 마디인 합집합--찾기로 계산한다. 합칠 때는 큰 뿌리를 작은 뿌리 밑에 붙이고, 찾기는 경로를 반감한다. 간선이 많으면 간선들을 크기로만 정한 범위로 나누어 여러 스레드가 동시에 합치는데, 어떤 순서로 합치든 같은 분할에 같은 대표가 남으므로 결과가 스레드 수와 무관하다. 14장에서 본 결정성의 원칙이 자료구조의 수준에서 지켜지는 예이다.
\end{instardust}

\section{맺음}

이 장의 세 자료구조는 모두 수십 년 된 고전이지만, 오늘날의 구현은 교과서의 판과 꽤 다르다. 힙은 캐시 줄에 맞추어 자식을 넷으로 늘리고, 해시 테이블은 분기 대신 SIMD 비교로 칸 열여섯 개를 한꺼번에 보며, 합집합--찾기는 결정성을 위해 대표를 고르는 규칙을 바꾼다. 점근적 복잡도는 같아도 기계의 구조와 결과의 재현성이라는 요구가 설계를 바꾸는 것이다. 다음 장에서는 이 책에서 가장 기본적인 자료구조인 다차원 배열과 데이터프레임의 열이 메모리에 어떻게 놓이는지를 다룬다.

\chapterreferences
